
\section{Evaluación y resultados}

\subsection{Metodología de medición (latencia E2E, confirmación, tasa de fallos)}

La evaluación se estructura con métricas cuantitativas de desempeño y consistencia funcional. Se miden tiempos de envío y confirmación de transacciones, latencia extremo a extremo desde interfaz, y proporción de fallos por tipo de operación.

\subsection{Resultados del prototipo (Testnet)}

En esta subsección se consolidan resultados de ejecución sobre testnet para los flujos de emisión, listado, compra, transferencia y redención. La presentación de resultados incluye tablas comparativas y análisis por percentiles cuando aplique.

\subsection{Evaluación antifraude (doble redención, reventa múltiple, QR estático)}

Se prueban escenarios de riesgo priorizados para validar que el sistema rechaza usos inválidos y conserva trazabilidad verificable del ticket a lo largo de su ciclo de vida.

\subsection{Discusión comparativa (contra estado del arte)}

Los resultados obtenidos se contrastan con referencias del estado del arte para discutir alcances, limitaciones y viabilidad técnica del enfoque basado en contratos inteligentes.

